% =============================================================================
% Chapter: System Evaluation. Defines each instrument briefly (the background
% introduces them) and reports its number. Figures are real, drawn from
% Report/metrics/{quantitative_results,ml_test_score}.md.
% =============================================================================

\chapter{Results: System Evaluation}
\label{ch:results-system}

This chapter evaluates the system. The forecast-accuracy results are in
Chapter~\ref{ch:results-models}. Following the single-case study guidelines of
Runeson and H\"ost~\cite{runeson2009casestudy}, the system is evaluated from its
source code, version-control and CI history, the ClickHouse experiment tracker,
and short targeted experiments. The figures are indicative of the system over
the study window rather than claims about long-run production operation, a
framing that matches comparable MLOps evaluations in the electricity-forecasting
domain~\cite{subramanya2022devops2mlops}. No controlled comparison against a
second platform is attempted.

The evaluation composes four established instruments introduced in
\S\ref{sec:bg-eval}, three objective measurements and one structured
self-assessment. An objective measurement is a number a script recomputes on
demand from code, git history, or the database, via
\texttt{scripts/report/extract\_metrics.py} and \texttt{reproduce\_run.py}. The
self-assessment is the ML Test Score, a rubric whose every item carries one line
of evidence a reader can check or contest. Between them the four instruments
evidence eight of Zarour et al.'s nine practices~\cite{zarour2025}. The ninth,
collaboration and communication, is qualitative and is documented in the system
chapter (Chapter~\ref{ch:system-implementation}) rather than measured here. Each
instrument is applied below and its result reported.

% -----------------------------------------------------------------------------
\section{ML Test Score}
\label{sec:res-mlts}

The ML Test Score rubric of Breck et al.~\cite{breck2017mltestscore} scores 28
tests across four sections on a $\{0, 0.5, 1\}$ scale, the final score being the
minimum of the four section sums so that the weakest area bounds the result.
Table~\ref{tab:res-mlts} reports the per-section sums. The full per-item evidence
is in Appendix~\ref{app:mlts-worksheet}.

\begin{table}[h]
  \centering
  \caption{ML Test Score self-assessment (min-of-sections rule). Per-item evidence in Appendix~\ref{app:mlts-worksheet}.}
  \label{tab:res-mlts}
  \renewcommand{\arraystretch}{1.3}
  \begin{tabularx}{\linewidth}{@{}>{\raggedright\arraybackslash}X r@{}}
    \toprule
    \textbf{Section (max 7)} & \textbf{Sum} \\
    \midrule
    Data tests                 & 2.5 \\
    Model development tests     & 3.0 \\
    ML infrastructure tests     & 4.0 \\
    Monitoring tests            & 4.0 \\
    \midrule
    \textbf{ML Test Score (min)} & \textbf{2.5} \\
    \bottomrule
  \end{tabularx}
\end{table}

The final score is \textbf{2.5}, bound by the data-testing section. On Breck et
al.'s interpretation scale this falls in the $(2,3]$ band, a system that is
reasonably tested but could automate more of its tests and procedures. Within
the data section the system scores for enforcing strict input schemas
(\texttt{enforce\_schema}) and for a tested feature-transformer pipeline, but
does not formally document the
feature-benefit, feature-cost, or privacy analyses the rubric also rewards, and
has no data-quality CI gate. We read this as the system's primary remaining
investment area. The score is deliberately
conservative. Two data items could each be argued to 0.5 (that the Lasso performs
feature selection, and that the feature set is grounded in Lago
et al.~\cite{lago2021}), which would lift the data section to 3.0 and the
headline to 3.0. We leave them at zero absent a documented analysis, so the
reported score is a lower bound.

% -----------------------------------------------------------------------------
\section{Reproducibility}
\label{sec:res-repro}

Reproducibility is reported in two senses, namely provenance coverage across all
runs and a worked re-run of one. The framing follows Gundersen and
Kjensmo~\cite{gundersen2018reproducibility}.

\begin{table}[h]
  \centering
  \caption{Provenance triple coverage across 15 live runs in the window.}
  \label{tab:res-repro-coverage}
  \renewcommand{\arraystretch}{1.3}
  \begin{tabularx}{\linewidth}{@{}>{\raggedright\arraybackslash}X r@{}}
    \toprule
    \textbf{Provenance component} & \textbf{Coverage} \\
    \midrule
    \texttt{git\_sha}                          & 66.7\,\% (10 / 15) \\
    \texttt{data\_start\_utc}                  & 100.0\,\% (15 / 15) \\
    \texttt{data\_end\_utc}                    & 100.0\,\% (15 / 15) \\
    \texttt{model\_artifact\_id}               & 100.0\,\% (15 / 15) \\
    \midrule
    \textbf{Full triple (all four)}             & \textbf{66.7\,\% (10 / 15)} \\
    \bottomrule
  \end{tabularx}
\end{table}

Of the 15 live runs in the window, 66.7\,\% carry the complete provenance triple
(Table~\ref{tab:res-repro-coverage}). The shortfall is entirely on
\texttt{git\_sha}: the data-window and artefact columns are populated for every
run from day one, while five early runs predate the CD pipeline that bakes
\texttt{GIT\_SHA} into the production image. From the first cd-deployed run
onward, all four components are present.

\begin{table}[h]
  \centering
  \caption{Reproducibility re-run: regenerated from logged provenance, compared against the stored prediction.}
  \label{tab:res-repro-rerun}
  \renewcommand{\arraystretch}{1.3}
  \begin{tabularx}{\linewidth}{@{}l >{\raggedright\arraybackslash}X@{}}
    \toprule
    \textbf{Field} & \textbf{Value} \\
    \midrule
    Config                       & \texttt{dk1\_lear} \\
    Target date                  & 2026-03-22 \\
    Delivery periods             & 24 \\
    Max absolute difference      & $7 \times 10^{-6}$\,EUR/MWh \\
    Verdict                      & numerically identical \\
    Command                      & \texttt{uv run python -m scripts.report.reproduce\_run --config dk1\_lear --date 2026-03-22} \\
    \bottomrule
  \end{tabularx}
\end{table}

For the demonstration (Table~\ref{tab:res-repro-rerun}), the only inputs supplied
were the run's committed config and target date, and the stored prediction was
loaded from \texttt{algo.experiment\_predictions}. For the deterministic LEAR
configuration on 2026-03-22, the maximum absolute difference across all 24
delivery periods was $7\times10^{-6}$\,EUR/MWh, five orders of magnitude below
any economically meaningful price unit, i.e.\ numerically identical up to
solver-level floating-point tolerance. LEAR is used for the headline
demonstration because it is deterministic given a fixed data window. The DNN and
Transformer carry seed and ensemble stochasticity whose controlled re-run is
left to future work.

% -----------------------------------------------------------------------------
\section{Drift detection as monitoring, not a retraining trigger}
\label{sec:res-drift}

The system refits every model from scratch daily on a rolling window at a cost of
seconds (Chapter~\ref{ch:system-implementation}, P6), so it is never stale between
refits and drift-\emph{triggered} retraining would add no freshness benefit,
following the cost-aware argument of Mahadevan and
Mathioudakis~\cite{mahadevan2024costaware}. Drift detection is therefore kept as a
\emph{monitoring and alerting} signal rather than a retraining trigger, and
\texttt{retrain\_on\_drift} defaults to false. Over the window the detector
flagged drift on 49 of 49 reports, an uninformatively high rate that calls for
recalibration.

% -----------------------------------------------------------------------------
\section{Maturity placement}
\label{sec:res-maturity}

The platform is positioned against the Google three-level maturity
model~\cite{google_mlops_maturity}. Of the five models reviewed by Zarour et
al.~\cite{zarour2025}, Google is chosen because its three levels match the
granularity of a single-host MVP, because its top level of CI/CD-pipeline
automation is plausibly attainable where Microsoft's Level 4 of full automated
operations would require multi-environment claims the platform does not make, and
because it is the most-cited maturity reference in that review, giving the
assignment an external touchpoint. Table~\ref{tab:maturity} assigns a status and
a Google level to each of Zarour's nine practices, with one line of evidence
apiece. The three status labels map to the evidence already presented. A practice
is \emph{Realised} when a concrete, production-running artefact closes it
end-to-end, and \emph{Minimal} when the mechanism exists but its scope is
intentionally restricted. Eight of the nine are Realised. The ninth, the feedback
loop, is Minimal, since drift-triggered retraining exists but is scoped to two
areas behind an opt-in flag, consistent with the monitoring-not-trigger argument
of \S\ref{sec:res-drift}. The levels themselves are descriptive ordinals and
support positioning rather than measurement. The quantitative load is carried by
the ML Test Score and reproducibility results above.

\begin{table}[h]
  \centering
  \caption{Maturity assessment of Zarour's nine practices, with the Google MLOps maturity model level \cite{google_mlops_maturity} assigned to each.}
  \label{tab:maturity}
  \renewcommand{\arraystretch}{1.3}
  \begin{tabularx}{\linewidth}{@{}l l c >{\raggedright\arraybackslash}X@{}}
    \toprule
    \textbf{Practice} & \textbf{Status} & \textbf{Google level} & \textbf{One-line evidence} \\
    \midrule
    P1 Automating CI/CD              & Realised & 2 & \texttt{pr\_test}/\texttt{pr\_ruff-lint}/\texttt{pr\_smoke} PR gates, \texttt{cd\_deploy} on every merge to \texttt{main}, 20 production deploys over the window at about 1.56 per week \\
    P2 Orchestrating workflows       & Realised & 2 & Prefect flows with per-area deployments on a Docker work pool \\
    P3 Ensuring reproducibility      & Realised & 2 & \texttt{GIT\_SHA} + data-window columns per run, containerised runtime in \texttt{docker-compose.yml} \\
    P4 Versioning                    & Realised & 2 & Git: \texttt{configs/}, numbered migrations, dashboards and alert rules all version-controlled \\
    P5 Collaboration \& communication & Realised & 1 & \texttt{CLAUDE.md}, \texttt{REPO\_OVERVIEW.md}, \texttt{RUNBOOK.md}, \texttt{CONTRIBUTING.md}, plus Slack alerts \\
    P6 Continuously train \& evaluate & Realised & 1 & Six \texttt{daily-forecast} deployments refit at 06:00 UTC, reconciliation scores live forecasts \\
    P7 Track \& log ML metadata      & Realised & 1 & \texttt{algo.experiment\_runs} + predictions + artifacts, plus the W\&B research surface \\
    P8 Continuous monitoring         & Realised & 2 & Loki plus three Grafana dashboards (Data Quality, Operational, System Health) and four file-based alert rules, with model metrics in W\&B \texttt{algo-prod} \\
    P9 Implementing feedback loops   & Minimal  & 1 & Drift-triggered retrain behind a config flag, two areas only \\
    \bottomrule
  \end{tabularx}
\end{table}

% =============================================================================
% Known gaps moved to the Conclusion (Future work) and the Discussion
% (Limitations) to avoid triplication.
% =============================================================================
